\chapter{Closing Remarks and Open Problems}\label{ch:closing}

\section{Uncountable sets and an infinite linear order game}

When does the converse of Proposition~\ref{prop:arrow} hold?
By the results of \ref{sec:cantor_game} such orderings must be $\omega_1$-free and $\omega_1^*$-free.
Separable and even ccc linear orderings have this property.
The separable case (the reals) was established in \cite{brian-clontz}.
Then, the next step is to study dense linear orders that are $\omega_1$-free, $\omega_1^*$-free but not separable.
A classical example of such ordering is the Aronszajn ordering.

\begin{definition}
    An \emph{Aronszajn line (A-line)} is an uncountable linearly ordered set that is $\omega_1$-free, $\omega_1^*$-free and contains no uncountable subset of the reals.
\end{definition}

Following \cite{A-orderings}, define a proper decomposition for A-lines.
Let $A$ be a dense Aronszajn line.
Let $D_0:=\emptyset$.
Given a countable $D_\alpha\subseteq A$, consider the equivalence relation on $A\backslash D_\alpha$: $x\sim_\alpha y$ if and only if there is no $z\in D_\alpha$ such that $x<z<y$.
Each $\sim_\alpha$ class is convex.
Let $T_\alpha$ be the set of all $\sim_\alpha$ classes and let $D_\alpha'$ be a selector of the $\sim_\alpha$ classes.
Since $D_\alpha$ is countable and dense in $D_{\alpha+1}:=D_\alpha\cup D_\alpha'$, then $D_{\alpha+1}$ must be countable because $A$ does not contain uncountable separable subsets.
Given $D_\alpha$ for $\alpha<\gamma$ and $\gamma$ limit, let $D_\gamma=\bigcup_{\alpha<\gamma}D_\alpha$.
Since $A$ does not contain copies of $\omega_1$ and $\omega_1^*$, $A=\bigcup_{\alpha<\omega_1}D_\alpha$.
Also, $(T,\supseteq)$ is an Aronszajn tree where $T=\bigcup_{\alpha<\omega_1}T_\alpha$.

\begin{question}
    Is there a dense Aronszajn line $A$ where Player II has a winning strategy in $BG(A,S)$ for all $S\subseteq A$?
\end{question}

Another class of lines where a similar decomposition can be performed is the class of nowhere separable Souslin continua.

\begin{definition}
    A \emph{Souslin continuum} is a non-separable dense, complete linearly ordered set that has the countable chain condition (ccc), i.e., every family of pairwise disjoint open intervals is countable.
    A Souslin continuum is \emph{nowhere separable} if no non-empty open interval is separable.
\end{definition}

Define a \emph{proper decomposition} for the nowhere separable Souslin continuum $L$ as follows.
Let $D_0:=D\neq\emptyset$ be any countable set.
Given a countable $D_\alpha\subseteq L$ consider the equivalence relation on $L\backslash\overline{D_\alpha}$:
$x\sim_\alpha y$ if and only if there is no $z\in D_\alpha$ such that $x<z<y$.
Again, each $\sim_\alpha$ class is convex.
Let $T_\alpha$ be the set of all $\sim_\alpha$ classes and let $D_\alpha'$ be a selector of the $\sim_\alpha$ classes.
Since $L$ has the ccc, then $T_\alpha$ is countable.
Let $D_{\alpha+1}=D_\alpha\cup D_\alpha'$.
For $\gamma<\omega_1$ limit let $D_\gamma=\bigcup_{\alpha<\gamma}D_\alpha$.

\begin{lemma}
    If $\{D_\alpha:\alpha<\omega_1\}$ is a proper decomposition of a nowhere separable Souslin continuum $L$, then $$L=\bigcup_{\alpha<\omega_1}\overline{D_\alpha}.$$
\end{lemma}

\begin{proof}
    Suppose that $x\in L$ and $x\notin \bigcup_{\alpha<\omega_1}\overline{D_\alpha}$.
    For each $\xi<\omega_1$ let $I_\xi^0:=[x]_\xi$ (the $\sim_\xi$ class of $x$).
    Note that $I_\eta^0\subseteq I_\xi^0$ for all $\xi<\eta$.
    \begin{claim}
        For all $\alpha<\omega_1$, $|T_\alpha|>1$.
    \end{claim}
    \begin{claimproof}
        Let $x_0\in D_\alpha$.
        Since $L$ is nowhere separable, $\overline{D_\alpha}$ is nowhere-dense so pick $y_1\in (L\backslash\overline{D_\alpha})\cap (-\infty, x_0)$ and $y_2\in (L\backslash\overline{D_\alpha})\cap (x_0,\infty)$.
        Then, $[y_1]_\alpha,[y_2]_\alpha\in T_\alpha$.
    \end{claimproof}
    For each $\xi<\omega_1$ let $I_{\xi+1}^1\in T_{\xi+1}$ such that $I_{\xi+1}^1\cap I_{\xi+1}^0=\emptyset$.
    Then, $\{I_{\xi+1}^1:\xi<\omega_1\}$ is an uncountable pairwise disjoint family of open sets, contradicting ccc.
\end{proof}

Using a proper decomposition to show the following remains an open question of great interest to the authors of \cite{MatosWiederholdSalvetti2025}.

\begin{conjecture}
    Let $L$ be a nowhere separable Souslin continuum.
    Then Player II has a winning strategy for $BG(L,S)$ if and only if $S$ is countable.
\end{conjecture}

\section{No Countable Basis for Borel Directed Graphs of Dichromatic Number at Least Three}\label{sec:dichrom-questions}
 
This section collects several natural questions suggested by the results presented in Section~\ref{sec:dichromatic}, the work of Raghavan and Xiao~\cite{RaghavanXiao2024}, and suggestions of Thornton.
 
\begin{question}[Minimality of the Raghavan--Xiao basis.]
	Raghavan and Xiao~\cite{RaghavanXiao2024} produce a continuum-size family of directed graphs that serves as a basis for Borel directed graphs with uncountable Borel dichromatic number.
	This family consists of pairwise Borel-incomparable directed graphs.
	\emph{Does every basis for this class have size~$\mathfrak{c}$?}
	Raghavan has noted (personal communication) that the directed graphs in his family are pairwise non-homomorphic, but it is unknown whether a smaller basis exists.
\end{question}

\begin{question}[Structural restrictions.]
	The result presented here applies to the full class of Borel directed graphs.
	\emph{Does a countable basis exist when one restricts to structurally simpler classes, say locally finite acyclic Borel directed graphs, or directed graphs generated by a single Borel function?}
	In the undirected case, Miller~\cite{Miller_2008_MeasurableChromaticNumbers} proved that there is no countable $\leq_B$-basis for graphs of the form~$G_f$ (where $f$ is a Borel function) with $\chi_B(G_f)\geq 3$.
	The analogous question for directed graphs generated by a single Borel function and dichromatic number remains open.
\end{question}	

\begin{question}[Borel dichromatic number and Borel order dimension.]
	Raghavan and Xiao~\cite{RaghavanXiao2024} established a close connection between Borel dichromatic number and a notion of Borel order dimension for Borel quasi-orders.
	\emph{Does the complexity result of Proposition~\ref{prop:main} yield new complexity results for Borel order dimension?}
\end{question}

\begin{question}[Dichotomies for complete multipartite graphs.]
	It is known that checking whether a Borel graph in which every connected component is complete $k$-partite admits a Borel $k$-coloring is~$\Pi^1_1$, but no satisfying dichotomy theorem is known for this problem.
	Thornton has suggested (personal communication) that a finite basis may exist for each~$k$, though its size appears to grow faster than~$n!$.
	\emph{Is there an explicit finite basis for each~$k$?  What is its growth rate?}
\end{question}

\begin{question}[Bounded width CSPs.]
	Among the constraint satisfaction problems studied in the Borel setting, the \emph{bounded width} problems occupy a distinguished position.
	Thornton~\cite{thornton2022algebraic} showed that the structures whose every solvable Borel instance has a Borel solution are exactly the width~$1$ structures, and proved partial complexity results for certain bounded width structures.
	Recent work of Grebík and Vidnyánszky~\cite{GV2025} shows that the split between easy and hard problems lies at a different place in the Borel setting than in the classical CSP dichotomy: in particular, there is no dichotomy for any Borel CSP that is not of bounded width.
	However, sharp complexity bounds within bounded width remain known only in very special cases.
	\emph{Can the exact complexity dividing line within bounded width CSPs be determined?}
	This problem appears to be very difficult; even identifying good illustrative examples of bounded width CSPs beyond the well-studied special cases remains open.
\end{question}

\begin{question}[Measurable arboricity.]
	In the undirected setting, the \emph{arboricity} of a graph (the minimum number of acyclic sets needed to cover the vertex set) is the undirected analogue of the dichromatic number.
	Several basic questions about arboricity in measurable combinatorics remain open.
	For instance:
	\emph{Can the $\mu$-measurable arboricity of a probability-measure-preserving graph differ from its (classical) arboricity by more than~$1$?}
	And: \emph{What is the $\mu$-measurable arboricity of the $k$-th power of the Bernoulli shift graph of~$F_2$?}
\end{question}